\chapter{Concorrência revisitada}
\label{CH:LOCK2}

Obter simultaneamente um bom desempenho paralelo, correção apesar da concorrência e um código compreensível é um grande desafio no design de kernels.
O uso direto de locks é o melhor caminho para garantir correção, mas nem sempre é possível.
Este capítulo destaca exemplos em que o xv6 é forçado a usar locks de forma elaborada e exemplos onde o xv6 utiliza técnicas semelhantes a locks, mas sem usá-los diretamente.
\section{Padrões de Bloqueio (Locking Patterns)}

Itens em cache frequentemente apresentam desafios para serem bloqueados.
Por exemplo, cópias de até {\tt NBUF} blocos de disco.
É crucial que um determinado bloco de disco tenha, no máximo, uma única cópia no cache; caso contrário, diferentes processos podem realizar alterações conflitantes em cópias diferentes do que deveria ser o mesmo bloco.
Cada bloco armazenado em cache é mantido em uma estrutura {\tt struct buf} \lineref{kernel/buf.h}{1}.
Uma {\tt struct buf} possui um campo de bloqueio (lock) que ajuda a garantir que apenas um processo use um determinado bloco de disco por vez.

No entanto, esse bloqueio não é suficiente: o que acontece se um bloco não estiver presente no cache e dois processos desejarem usá-lo ao mesmo tempo?
Nesse caso, não há uma {\tt struct buf} (já que o bloco ainda não foi armazenado em cache), e, portanto, não há nada para bloquear.

O xv6 resolve essa situação associando um bloqueio adicional ({\tt bcache.lock}) ao conjunto de identidades dos blocos armazenados em cache.
O código que precisa verificar se um bloco está no conjunto de blocos armazenados em cache deve segurar o bloqueio {\tt bcache.lock}; depois que o código encontra o bloco e a {\tt struct buf} necessária, ele pode liberar o {\tt bcache.lock} e bloquear apenas o bloco específico.
Esse é um padrão comum: um bloqueio para o conjunto de itens e um bloqueio por item.

Normalmente, a mesma função que adquire um bloqueio também o libera. No entanto, uma forma mais precisa de enxergar isso é considerar que um bloqueio é adquirido no início de uma sequência que deve parecer atômica e é liberado quando essa sequência termina.
Se a sequência começa e termina em funções diferentes, ou em threads diferentes, ou até mesmo em CPUs diferentes, então a aquisição e a liberação do bloqueio devem acompanhar esse comportamento.
A função do bloqueio é forçar outros usos a aguardarem, e não vincular um pedaço de dados a um agente específico.
Um exemplo disso é a função {\tt acquire} em {\tt yield}, que pode ser chamada em um thread diferente do processo que originalmente adquiriu o bloqueio.
Outro exemplo ocorre quando o código frequentemente entra em estado de espera enquanto lê o disco; ele pode despertar em uma CPU diferente, o que significa que o bloqueio pode ser adquirido e liberado em CPUs distintas.

Liberar um objeto que é protegido por um bloqueio embutido no próprio objeto é uma tarefa delicada, pois possuir o bloqueio não é suficiente para garantir que a liberação será correta.
O problema ocorre quando algum outro thread está aguardando em {\tt acquire} para usar o objeto; liberar o objeto implicitamente libera o bloqueio embutido, o que causará mau funcionamento no thread que estava aguardando.
Uma solução para esse problema é rastrear quantas referências ao objeto existem, garantindo que ele só será liberado quando a última referência desaparecer.
Um exemplo disso pode ser observado em {\tt pipeclose}, onde {\tt pi->readopen} e {\tt pi->writeopen} rastreiam se há descritores de arquivos referindo-se ao pipe.

Um motivo pelo qual condições de corrida podem quebrar programas é que, se não houver locks ou construções equivalentes, o compilador pode gerar código de máquina que lê e grava na memória de maneiras bem diferentes do código C original. Por exemplo, o código de máquina de uma thread que chama
\texttt{killed} poderia copiar \lstinline{p->killed} para um registrador e ler apenas esse valor em cache; isso significa que a thread talvez nunca veja nenhuma escrita em
\lstinline{p->killed}. Os locks evitam esse tipo de caching.

% sleep locks
% hand-over-hand locking em namei
% Uso de refcount pelo namei para evitar modificações durante o acesso,
% e lock quando efetivamente em uso.
% Exemplo onde deadlock é evitado? namei?
% Criar uma thread para contornar um problema de ordem de locks?

\section{Padrões semelhantes a locks}

Em muitos pontos, o xv6 utiliza contadores de referência ou flags de forma semelhante a locks para indicar que um objeto está alocado e não deve ser liberado ou reutilizado. O campo {\tt p->state} de um processo age dessa forma, assim como os contadores de referência em estruturas {\tt file}, {\tt inode} e {\tt buf}. Embora em cada caso um lock proteja a flag ou o contador de referência, é este último que impede que o objeto seja liberado prematuramente.

O sistema de arquivos utiliza os contadores de referência de {\tt struct inode} como uma espécie de lock compartilhado que pode ser mantido por múltiplos processos, a fim de evitar deadlocks que ocorreriam se o código utilizasse locks comuns. Para o diretório nomeado por cada componente do caminho, um de cada vez. No entanto, {\tt namex} deve liberar cada lock ao final do laço, pois se mantivesse múltiplos locks poderia entrar em deadlock consigo mesmo se o caminho incluísse um ponto (por exemplo, {\tt a/./b}). Também poderia entrar em deadlock com uma busca concorrente envolvendo o diretório e {\tt ..}. Como explicado no Capítulo~\ref{CH:FS}, a solução é o laço carregar o inode do diretório para a próxima iteração com seu contador de referência incrementado, mas sem estar bloqueado (locked).

Alguns itens de dados são protegidos por mecanismos diferentes em momentos diferentes, e às vezes podem estar protegidos contra acesso concorrente de forma implícita, pela estrutura do código do xv6, em vez de por locks explícitos. Por exemplo, quando uma página física está livre, ela é protegida por {\tt kmem.lock}, e quando está alocada para uma estrutura {\tt proc}, é protegida por um lock diferente (o lock embutido {\tt pi->lock}). Se a página for realocada para a memória de usuário de um novo processo, ela não é protegida por nenhum lock. Em vez disso, o fato de o alocador não fornecer essa página para nenhum outro processo (até que ela seja liberada) é o que a protege contra acesso concorrente.

A propriedade da memória de um novo processo é complexa: primeiro o pai a aloca e a manipula em {\tt fork}, depois o filho a utiliza, e (após o término do filho) o pai volta a ser o proprietário da memória e a libera com {\tt kfree}. Há duas lições aqui: um objeto de dados pode ser protegido contra concorrência de maneiras diferentes em diferentes momentos da sua vida útil, e essa proteção pode assumir a forma de estrutura implícita em vez de locks explícitos.

Um exemplo final semelhante a um lock é a necessidade de desabilitar interrupções ao redor de chamadas de interrupção, o que torna o código que as chama atômico em relação a interrupções de temporizador que poderiam forçar uma troca de contexto, e assim mover o processo para uma CPU diferente.

\section{Sem locks}

Existem alguns pontos em que o xv6 compartilha dados mutáveis sem usar nenhum tipo de lock. Um exemplo é na implementação de spinlocks, embora se possa considerar que as instruções atômicas do RISC-V dependem de locks implementados em hardware. Outro exemplo é a variável {\tt started} em {\tt main.c}, usada até que a CPU zero termine de inicializar o xv6; o uso de {\tt volatile} garante que o compilador realmente gere instruções de leitura (load) e escrita (store).

O xv6 contém casos em que uma CPU ou thread escreve alguns dados, e outra CPU ou thread lê esses dados, mas não há um lock específico dedicado à proteção desses dados. Por exemplo, em {\tt fork}, o processo pai escreve as páginas de memória do usuário do processo filho, e o filho (um thread diferente, talvez em uma CPU diferente) lê essas páginas; nenhum lock protege explicitamente essas páginas. Isso não é exatamente um problema de lock, já que o filho não começa a executar até que o pai tenha terminado de escrever. No entanto, é um possível problema de ordenação de memória (veja o Capítulo~\ref{CH:LOCK}), pois, sem uma barreira de memória, não há garantia de que uma CPU verá as escritas feitas por outra CPU. Entretanto, como o pai libera locks e o filho adquire locks ao iniciar, as barreiras de memória presentes em {\tt acquire} e {\tt release} garantem que a CPU do filho veja as escritas feitas pelo pai.
\section{Paralelismo}

O uso de travas (locks) trata-se principalmente de suprimir o paralelismo em prol da correção. Como o desempenho também é importante, os desenvolvedores do kernel frequentemente precisam pensar em como usar travas de maneira que se obtenha tanto a correção quanto o paralelismo. Embora o xv6 não seja projetado sistematicamente para alto desempenho, ainda vale a pena considerar quais operações do xv6 podem ser executadas em paralelo e quais podem entrar em conflito devido a travas.

Os pipes no xv6 são um exemplo de paralelismo relativamente bom. Cada pipe possui sua própria trava, o que permite que processos diferentes possam ler e escrever em pipes diferentes em paralelo, em diferentes CPUs. Para um dado pipe, no entanto, o escritor e o leitor precisam esperar um ao outro liberar a trava; eles não podem ler/escrever no mesmo pipe ao mesmo tempo. Também é verdade que uma leitura de um pipe vazio (ou uma escrita em um pipe cheio) deve bloquear, mas isso não ocorre devido ao esquema de travas.

O escalonamento de contexto (context switching) é um exemplo mais complexo. Duas threads do kernel, cada uma executando em sua própria CPU, podem chamar {\tt yield}, {\tt sched} e {\tt swtch} ao mesmo tempo, e essas chamadas serão executadas em paralelo. As threads mantêm uma trava, mas são travas diferentes, então elas não precisam esperar uma pela outra. No entanto, uma vez dentro da função {\tt scheduler}, as duas CPUs podem entrar em conflito por travas ao buscar na tabela de processos por um processo que esteja com o estado {\tt RUNNABLE}. Ou seja, o xv6 provavelmente obtém um ganho de desempenho com múltiplas CPUs durante a troca de contexto, mas talvez não tanto quanto poderia.

Outro exemplo são chamadas concorrentes à função {\tt fork} feitas por diferentes processos em diferentes CPUs. As chamadas podem precisar esperar umas pelas outras por conta da {\tt pid_lock} e da {\tt kmem.lock}, bem como pelas travas por processo necessárias para buscar na tabela de processos por um processo com estado {\tt UNUSED}. Por outro lado, os dois processos que estão executando {\tt fork} podem copiar páginas de memória do usuário e formatar páginas da tabela de páginas totalmente em paralelo.

O esquema de travas em cada um dos exemplos acima sacrifica o desempenho paralelo em certos casos. Em todos os casos, é possível obter mais paralelismo usando um projeto mais elaborado. Se vale a pena ou não depende de vários detalhes: com que frequência as operações relevantes são invocadas, quanto tempo o código permanece com uma trava disputada em uso, quantas CPUs podem estar executando operações conflitantes ao mesmo tempo, se outras partes do código são gargalos mais restritivos. Pode ser difícil prever se um determinado esquema de travas causará problemas de desempenho, ou se um novo projeto será significativamente melhor, portanto medições com cargas de trabalho realistas são frequentemente necessárias.

\section{Exercícios}

\begin{enumerate}

\item Modifique a implementação de pipes do xv6 para permitir que uma leitura e
uma escrita no mesmo pipe ocorram em paralelo em CPUs diferentes.

\item Modifique a função \texttt{scheduler()} do xv6 para reduzir a contenção por travas
quando diferentes CPUs estiverem procurando por processos prontos para executar ao mesmo tempo.

\item Elimine parte da serialização na função \texttt{fork()} do xv6.

\end{enumerate}




